\documentclass[bibtotoc]{scrartcl}   %bibtotoc ist daf?r da, dass das Literaturverzeichnis auch im Inhaltsverzeichnis angezeigt wird, liststotoc ist f?r Abbildungsverzeichnis und Tabellenverzeichnis im Inhaltsverzeichnis
%\pagestyle{headings} %Kapitel?berschrift und Seitennummer in Kopfzeile (au?er dort, wo neues Kapitel beginnt



%Paket für ein deutsches Literaturverzeichnis
% \usepackage{natbib}
% \bibliographystyle{agsm}

%Paket für ein deutsches Literaturverzeichnis
\usepackage[authoryear]{natbib}

% \bibliographystyle{natdin}
\bibliographystyle{natdin}

\bibliographystyle{plainnat}

%\setcounter{secnumdepth}{4} %f?r die Auflistungstiefe im Inhaltsverzeichnis. So w?rde jetzt bis zur 4. Unterteilung alles aufgelistet.
%\setcounter{tocdepth}{4}
\usepackage[utf8x]{inputenc}
%\usepackage[applemac]{inputenc} %setzt Sonderzeichen um
\usepackage{lmodern} %f?r die Trennung der W?rter
\usepackage[ngerman]{babel} % f?r die deutsche Beschriftung von z.B. Grafiken
\usepackage[babel, german=quotes]{csquotes} %Erzeugung von Anführungszeichen durch \enquote{text}
\usepackage[T1]{fontenc}
\usepackage{longtable} % erm?glicht mehrseitige Tabellen
\usepackage{booktabs} % entzerrt die Tabellenzeilen und bietet verschiedene dicke Unterteilungslinien
\usepackage{graphicx} %um Grafiken einbinden zu k?nnen
% \usepackage{float} % unterst?tzt bei der Positionierung von Grafiken und Tabellen
\usepackage{color}
\definecolor{black}{gray}{0}% 10% gray
\usepackage[colorlinks=true, linkcolor=black, urlcolor=black, citecolor=black]{hyperref}  %verwandelt alle Kapitel?berschriften, Verweise aufs Literaturverzeichnis und andere Querverweise in PDF Hyperlinks und stellt sie angenehm unauff?llig dar
\usepackage[normalem]{ulem} % l?sst Unterstreichungen umbrechen
\usepackage{textcomp}
\usepackage{multirow}
\usepackage{pdfpages} %zum Einbinden von PDFs
\usepackage{url}
\usepackage{colortbl}
\usepackage{setspace} %um den Zeilenabstand zu regulieren
\usepackage{geometry}
\geometry{a4paper, left=3cm, right=2cm, top=2cm}
\usepackage{paralist} %braucht man für formatierte Aufzählungen. Mit \begin{compactitem} \item ... und \end{compactitem} bekommt man keine Leerzeichen zwischen den aufgezählten Punkten
%\usepackage{remreset}
\usepackage[multiple]{footmisc}



\onehalfspacing

%\hypersetup{linkcolor=black, urlcolor=black} %um URLs nicht rosa anzeigen zu lassen


%\usepackage{tocloft}

%\renewcommand{\cftfigpresnum}{Abb. }
%\renewcommand{\cftfigaftersnum}{:}

%% Befehl fur Literaturverweise mit Seitenangabe
%% Bsp: \site{Rupp2004}{213}
\newcommand{\site}[2]{
  \citep[S. #2]{#1}
}

% TODO Box
\usepackage{mdframed}
\newcommand{\todo}[1]{
  \vspace{0.2cm}
  \begin{mdframed}[skipabove=3pt, skipbelow=3pt, backgroundcolor=green]
    TODO: #1
  \end{mdframed}
  \vspace{0.2cm}
}

\date{\today}


\begin{document}

%!TEX root = ../hausarbeit.tex
\begin{titlepage}



%Logo der Fachhochschule Köln
\begin{figure}[!ht]
	\centering
		\includegraphics[natwidth=920pt, natheight=95pt, width=1.0\textwidth]{Bilder/LogoUni.png}


\end{figure}

\begin{center}

\vspace{0,8cm}

%vorgelegt an der...
\begin{large}
Universität zu Köln\\
%Humanwissenschaftliche Fakultät\\
Erziehungswissenschaft (1-Fach)\\
BM 3 - Pädagogik und Gesellschaft\\
\vspace{1,2cm}
\begin{scshape}
Hausarbeit\\ 
\end{scshape}
\end{large}


\vspace{1.2cm}

%Deutscher Titel
\begin{rmfamily}
\textbf{\huge Jugend und Medien}\\
\LARGE -\\
	Welchen Stellenwert hat Medienkompetenz in Deutschland?\\
\normalsize
\end{rmfamily}

\vspace{1.7cm}
%
%%Englischer Titel
%\begin{rmfamily}
%\textbf{\LARGE classic vs. agile project management}\\
%\large -\\
%	Compare klassic and agiel procedure model relating to their suitability in the media sector, especially in 3D projects
%\normalsize
%\end{rmfamily}

%\vspace{1.2cm}

\begin{large}
Seminar: Jugend und Medien\\ [0.8em]
\end{large}

%\vspace{1.0cm}

%Autor der Bachelorarbeit und die Prüfer
\begin{tabular}{rl}
        Seminarleiterin:  &  Ilona Cwielong\\
       			       &  \small Universität zu Köln \\[0.8em]
\end{tabular}
							  
\begin{large}
Wintersemester 2012/2013\\
\todo{Abgabedatum eingeben}
Abgabe der Arbeit: xx.xx.2013\\
\end{large}

\vspace{1.7cm}

%\end{center}

%\newpage
%\thispagestyle{empty}

%Kontaktmöglichkeiten des Autors und der Prüfer
%\begin{center}

%Autor der Bachelorarbeit und die Prüfer
\begin{tabular}{rl}
        eingereicht von:  &  Birgit Schlotter\\
       			    	  &  Humboldtstraße 15\\
			      		  &  50676 Köln\\
			      		  &  Matrikel-Nr.: 5587077\\
						  &  Tel.: 01 77 - 6 41 66 05\\
						  &  E-mail: bschlott@smail.uni-koeln.de\\
						  &  1. Semester\\[0.8em]
\end{tabular}



%\begin{tabular}{rl}
%						
%	eingereicht von: &   Birgit Schlotter\\
%			&  Humboldtstraße 15\\
%			&  50676 Köln\\
%			&  Tel.: 01 77 - 6 41 66 05\\
%			&  E-mail: bschlott@smail.uni-koeln.de\\
%			&  1. Semester\\[0.8em]
							%&  \quad 1. Semester\\[2.0em]
							
%\end{tabular}
							
			
							
							
\end{center}

\end{titlepage}

\clearpage

\singlespacing %Inhaltsverzeichnis nur mit 1-Zeiligem Abstand
\tableofcontents  %Inhaltsverzeichnis
\thispagestyle{empty}
\clearpage

\onehalfspacing

%!TEX root = ../hausarbeit.tex
\section{Einleitung}

Bla bla
%!TEX root = ../hausarbeit.tex
\section{Was ist Medienkompetenz}

Bla bla
%!TEX root = ../hausarbeit.tex
\section{Welche Medien nutzen Jugendliche}

Die Ausstattung mit den unterschiedlichsten Medien hat in den letzten Jahren sowohl in privaten Haushalten als auch in der Schule deutlich zugenommen. Meilensteine waren dabei der Computer/Laptop, das Smartphone, die Spielekonsole und natürlich das Internet. Die Nutzung der Medien ist für die meisten Jugendlichen heutzutage vollständig in den Alltag integriert. \todo{schreiben, dass es mehrere Studien gibt, ein paar aufzählen und dann sagen, dass man sich, um den Rahmen nicht zu sprengen, auf die Jim Studie begrenzt...}

Die JIM Studie\footnote{JIM (Jugend, Information, (Multi-)Media) ist eine Basisstudie zum Medienumgang 12- bis 19-Jähriger in Deutschland. Herausgeber ist der Medienpädagogische Forschungsverbund Südwest unter der Leitung von Peter Behrens und Thomas Rathgeb} bietet repräsentative und objektive Daten und Fakten rund um den aktuellen Medienumgang von 12- bis 19-Jährigen. Dafür wurde in der Zeit vom 7. Mai bis 17. Juni 2012 eine repräsentative Stichprobe von 1.201 Jugendliche im Alter von 12 bis 19 Jahren telefonisch befragt. Es nahmen 49\% Mächen und 51\% Jungen teil die zu 86\% Schüler waren und zu 54\% das Gymnasium besuchten. 

In \citet{jim12} werden folgende Daten zur akteullen Mediennutzung angegeben: 100\% der Jugendlichen gaben an, dass in ihrem Haushalt ein Laptop/Computer vorhanden ist. 79\% der Mädchen und 85\% der Jungen gaben an, einen eigenen Laptop/Computer zu besitzen. In 98\% der Fälle ist ein Handy, ein Fernseher und Internetzugang im Haushalt vorhanden, 98\% der Mädchen und 95\% der Jungen besitzen ein eigenes Handy (davon entfallen 47\% auf ein Smartphone), 55\% der Mädchen und 64\% der Jungen verfügen über einen eigenen Fernseher und 85\% der Mädchen und 88\% der Jungen können von ihrem Zimmer aus auf das Internet zugreifen. 69\% gaben an, dass sie feste und tragbare Spielekonsolen im eigenen Haushalt haben, davon besitzen 45\% der Mädchen und 57\% der Jungen eigene Konsolen, die sie nicht mit der Familie teilen müssen. 
%!TEX root = ../hausarbeit.tex
\section{Vorteile der Mediennutzung}

Bla bla
%!TEX root = ../hausarbeit.tex
\section{Gefahren der Mediennutzung}

Bla bla
%!TEX root = ../hausarbeit.tex
\section{Projekte zur Stärkung von Medienkompetenz}

Bla bla
%!TEX root = ../hausarbeit.tex
\section{Fazit}

bla bla





\pagebreak

%\listoffigures %Abbildungsverzeichnis
%\listoftables %Tabellenverzeichnis


% \appendix
\bibliography{literatur}

% The command \nocite{*} causes all items in the database to be included in the references, regardless of whether or not they are cited in the paper.
% \nocite{*}

\end{document}


